\chapter{Diskussion und Schlussbetrachtung}

Das Kapitel beantwortet die vier Forschungsfragen, grenzt die Aussagekraft der internen und
externen Ergebnisse ein und leitet daraus ein betriebliches Einsatz- und Pilotkonzept ab.


\section{Interpretation und Beantwortung der Forschungsfragen}

\textbf{FF1 -- task-adaptive Repräsentationsanpassung.}
R0 bildet bereits einen allgemein selbstüberwacht vortrainierten Ausgangspunkt. Der in dieser
Arbeit gemessene Effekt beschreibt daher den zusätzlichen Nutzen der task-adaptiven
Weiteranpassung. RU und RT verbessern die Downstream-Leistung gegenüber R0, während
RUT mit beiden Probes die höchste mittlere TPR erreicht. Der N10-Vergleich bestätigt diesen
Vorteil auf FINAL: Auf \nolinkurl{OFFICIAL_TEST} liegt RUT bei 20.000 Labels um
+6,08 Prozentpunkte über R0; in den vier zusätzlichen Evaluationsbedingungen beträgt der
Unterschied +5,51 bis +6,26 Prozentpunkte und bleibt nach Holm-Korrektur signifikant.

Damit zeigt sich ein zusätzlicher Repräsentationsnutzen der Anpassung auf dem ungelabelten
Webseitenpool, obwohl die Ausgangsencoder bereits selbstüberwacht vortrainiert sind.
Die in Kapitel 3 dargestellten spezialisierten URL-Repräsentationsverfahren verdeutlichen
die Relevanz einer domänennäheren Repräsentationsbildung für URL-basierte
Sicherheitsaufgaben (vgl. [Maneriker et al. (2021), PDF-S. 1--2]; [Wang et al. (2023), Abstract, o.\,S.]; [Li et al. (2025), S. 1--2 und 12]).
Die vorliegende Untersuchung betrachtet diesen Zusammenhang für die zusätzliche
task-adaptive Anpassung zweier Informationszweige innerhalb desselben Versuchsprotokolls.

Die faktorielle Betrachtung zeigt zugleich, dass sich die einzelnen Adaptationsgewinne nicht
vollständig addieren. Die Interaktion beträgt $-5{,}10$ Prozentpunkte für die lineare und
$-2{,}83$ Prozentpunkte für die MLP-Probe. URL- und HTML-Text-TAPT liefern damit jeweils
einen positiven Beitrag, weisen bei gemeinsamer Anpassung jedoch eine teilweise Überlappung
ihrer Effekte auf. RUT bleibt dennoch die Variante mit der höchsten absoluten
Downstream-Leistung.

Der vorgelagerte kontrastive Versuchsblock ist davon getrennt zu interpretieren. Die
getesteten kontrastiven Varianten lagen unter der damaligen Basis- und MLM-Konfiguration,
wurden jedoch nicht unter dem finalen R0/RU/RT/RUT-Protokoll erneut trainiert. Daraus
lässt sich keine allgemeine Aussage über kontrastives Lernen gegenüber MLM-basiertem
TAPT ableiten.


\textbf{FF2 -- Einfluss der Labelverfügbarkeit.}
Der RUT-Vorteil bleibt über alle untersuchten internen Labelbudgets von 2.000 bis
200.000 Instanzen bestehen. Bei kleineren Trainingsmengen ist der Abstand zu R0 besonders
ausgeprägt; mit zunehmender gelabelter Supervision verringert er sich, ohne vollständig zu
verschwinden. Die task-adaptive Repräsentationsanpassung verbessert damit nicht nur einen
einzelnen Betriebspunkt, sondern den Verlauf der gesamten Lernkurve.

Für die Einordnung ist die Trennung der beiden Lernphasen relevant. Der TAPT-Schritt
verwendet den 200.000 Instanzen umfassenden SSL-Pool ohne Downstream-Labels. Erst im
anschließenden Klassifikationstraining wird der Umfang der verfügbaren gelabelten Supervision
variiert. Die Sensitivität gegenüber reduziertem gelabeltem Trainingsumfang wird auch für
spezialisierte URL-Repräsentationen als relevante Evaluationsdimension betrachtet
(vgl. [Li et al. (2025), S. 12]). Die B95-Auswertung weist ausschließlich auf DEV auf eine frühere relative Sättigung von RUT hin; auf FINAL erreichen R0 und RUT das jeweilige relative Ziel bei denselben geprüften Budgets.
Da sich B95 jeweils auf das eigene 200k-Endniveau einer Variante bezieht, beantwortet diese
Kennzahl jedoch nicht, ob ein reduziertes RUT-Budget auch das absolute Niveau einer
vollständig trainierten Referenz erreicht. Diese Frage wird durch FF3 gesondert geprüft.

Die externe Evaluation auf dem Phishing Website Dataset von Putra (2023) erweitert diese
Einordnung über den PhreshPhish-Datenbestand hinaus. Bei identischen gelabelten
Zielinstanzen erreicht RUT über die vier Budgets von 200 bis 2.000 Webseiten eine um
2,54 Prozentpunkte höhere AP-AULC als R0 (95\,\%-KI [2,16; 2,91], 10/10 positive
Paare, $p=0{,}001953$). Auch die vier einzelnen Budgetvergleiche fallen mit
AP-Vorteilen zwischen +2,21 und +2,79 Prozentpunkten zugunsten von RUT aus und bleiben
nach Holm-Korrektur signifikant. Der Nutzen der task-adaptierten Repräsentation zeigt sich damit nicht ausschließlich innerhalb
der ursprünglichen PhreshPhish-Lernkurve. Im untersuchten externen Datenbestand bildet RUT
auch bei begrenzter neuer gelabelter Supervision die günstigere Ausgangsbasis für die
Downstream-Klassifikation. Der Befund stützt damit die Annahme, dass eine zuvor auf
ungelabelten Webseitendaten angepasste Repräsentation die spätere Anpassung an eine neue
Zielumgebung erleichtern kann.


Die ergänzende Putra-Zielanpassung verschlechtert dagegen die AP-AULC um
1,14 Prozentpunkte. Mehr selbstüberwachtes Vortraining führt folglich in diesem
Protokoll nicht zwangsläufig zu besseren Repräsentationen. Der Versuch isoliert
weder die Ursache der Verschlechterung noch eine optimale Adaptationsstärke;
insbesondere Lernrate, Trainingsdauer und Korpusgröße wurden hier nicht systematisch variiert.

\textbf{FF3 -- Erreichen des Referenzniveaus bei reduziertem Labelbudget.}
Die formale Budgetanalyse zeigt, dass 50.000 Labels die kleinste untersuchte RUT-Stufe
sind, die das festgelegte Vergleichsniveau gegenüber R0@200k über alle fünf
Evaluationsbedingungen einschließlich Holm-Korrektur erreicht. Die mittlere TPR von
RUT50k liegt dabei um +3,24 bis +3,62 Prozentpunkte über R0@200k; zugleich fällt die
realisierte FPR niedriger aus.

Auch gegenüber TF-IDF/XGBoost@200k ist 50.000 die kleinste statistisch gestützte
Budgetstufe. RUT20k liegt zwar bereits in allen fünf Bedingungen mindestens auf der
mittleren XGBoost-TPR und weist gleichzeitig eine niedrigere mittlere FPR auf, besteht
jedoch die Korrektur über das vollständige Budgetraster nicht. Mit RUT50k beträgt der
TPR-Abstand zu XGBoost je nach Evaluationsbedingung +4,59 bis +5,58 Prozentpunkte.
Da sich bei diesem Vergleich sowohl Repräsentation als auch Klassifikator ändern, dient
XGBoost als zusätzliche Referenz für das erreichbare Leistungsniveau und nicht als isolierter
Nachweis eines TAPT-Effekts.

Im untersuchten Linear-Probe-Protokoll reichen damit 50.000 gelabelte Instanzen aus, um
das festgelegte Leistungsniveau beider mit 200.000 Labels trainierten Referenzen zu erreichen.
Bezogen auf die Downstream-Supervision entspricht dies einer Reduktion um 75\,\%. Der
vorausgehende TAPT-Schritt auf dem vollständigen ungelabelten Webseitenpool bleibt davon
unberührt. Die Margensensitivität bestätigt die 50k-Entscheidung zusätzlich über den
untersuchten Bereich von 0 bis 2 Prozentpunkten. Die zusätzlichen internen Evaluationsbedingungen zeigen, dass diese Relation nicht auf den
vollständigen Finalbestand beschränkt ist. Dies ist insbesondere vor dem in Kapitel 3
diskutierten Problem datensatzabhängiger Evaluationsresultate relevant
(vgl. [Rashid et al. (2024), Abstract, o.\,S.]; [Dalton et al. (2026), S. 5--6]). Die Domain-,
Template- und Zeitbedingungen sind jedoch nicht als kausale Einzeltests verschiedener Formen
von Distribution Shift zu verstehen. Sie prüfen, ob die in FF1 und FF2 identifizierten
Leistungsrelationen unter veränderter Zusammensetzung der Testdaten erhalten bleiben.

Die realisierte FPR liegt insbesondere auf den domainbezogenen Teilmengen teilweise oberhalb
des auf CAL angestrebten Betriebspunkts. Eine stabile Rangfolge der Modelle bedeutet daher
nicht, dass ein einmal kalibrierter Schwellenwert unter jeder Datenverteilung dieselbe FPR
erzeugt. Repräsentationsrobustheit und Betriebspunktkalibration sind entsprechend als getrennte
Fragestellungen zu behandeln. Auch die ergänzende Prävalenzanalyse aus Abschnitt 6.6 betrifft eine andere Eigenschaft der
Evaluation. Unter konstant gehaltenem klassenspezifischem Fehlerverhalten sinkt
$P_{\geq 90}$ bei einer rechnerischen Phishing-Base-Rate von 1\,\% auf 80,74\,\%
für TRI Full beziehungsweise 84,27\,\% für die Kaskade. Die Analyse verdeutlicht damit
die Abhängigkeit precision-basierter Kennzahlen von der Klassenprävalenz.


\textbf{FF4 -- multimodale Systemkomponenten und URL-first-Verarbeitung.}
Die sequenzielle Komponentenablation zeigt positive zusätzliche Beiträge aller drei
Erweiterungsschritte. Den größten Zuwachs liefert HTML-Text mit +5,184 Prozentpunkten
TPR gegenüber URL-only. Der DOM-Zweig erhöht die TPR anschließend um weitere
+0,728 Prozentpunkte und das lernbare Gating um +0,576 Prozentpunkte. Alle drei
benachbarten Vergleiche bleiben nach Holm-Korrektur signifikant.

Die Ergebnisse zeigen damit für das implementierte System, dass die in Kapitel 3 diskutierten
textuellen und strukturellen Webseiteninformationen über die URL hinaus zusätzliche
Klassifikationsinformation bereitstellen können (vgl. [Aljofey et al. (2022), S. 1--2 und 13]; [Yoon et al. (2024), S. 13 und 15]). Gleichzeitig unterscheiden sich ihre marginalen Beiträge:
Der HTML-Text erzeugt den größten zusätzlichen Effekt, während DOM und Gating kleinere,
aber statistisch gestützte Verbesserungen liefern. TRI Full erreicht auf \nolinkurl{OFFICIAL_TEST} 96,12\,\% TPR bei 0,666\,\% FPR und
bildet die leistungsstärkste Variante des implementierten Gesamtsystems. Die
URL-first-Kaskade ergänzt diese Architektur um eine gestufte Verarbeitung. In der
eingefrorenen Routingkonfiguration werden 14,93\,\% der Webseiten an den vollständigen
TRI-Pfad eskaliert; die TPR beträgt 95,77\,\%. Gleichzeitig sinkt die Median-Latenz von
30,18 auf 11,95\,ms und der gemessene In-Memory-Durchsatz steigt von 183,2 auf
447,0 Seiten pro Sekunde. Die Kaskade realisiert damit das in Kapitel 3 aufgegriffene
Prinzip einer mehrstufigen Verarbeitung, ohne den multimodalen Pfad für unsichere Fälle
aufzugeben (vgl. [Liu et al. (2021), Abstract, o.\,S.]). Der Frontend-Vergleich verbindet schließlich Repräsentations- und Systemebene. Bei derselben
auf Development festgelegten TPR-Toleranz von einem Prozentpunkt sinkt die mittlere
Eskalationsrate auf \nolinkurl{OFFICIAL_TEST} von 35,746\,\% mit URL\_BASE
auf 23,514\,\% mit URL\_TAPT (fünf Läufe).
Die task-adaptierte URL-Repräsentation ermöglicht innerhalb dieses Routingprotokolls damit
häufiger eine Entscheidung in der ersten Verarbeitungsstufe. Der Repräsentationsvorteil wirkt
sich folglich nicht nur auf die Klassifikationsleistung eines einzelnen Modells, sondern auch
auf die Verarbeitungstiefe des Gesamtsystems aus.

Die zentrale Forschungsfrage wird damit auf drei Ebenen beantwortet: zusätzlicher
Repräsentationsnutzen durch TAPT, geringerer gelabelter Trainingsbedarf im Probe-Protokoll
und eine steuerbare Verarbeitungstiefe des darauf aufbauenden multimodalen Systems.


\section{Limitationen und Aussagegrenzen}

Die internen Generalisierungsanalysen beruhen auf demselben PhreshPhish-Data-Freeze.
Der ergänzende Integritätsaudit schließt für die prüfbaren OOD-Instanzen exakte
Überschneidungen der Domains beziehungsweise Template-Hashes mit SUP, DEV, CAL und
dem SSL/TAPT-Pool aus. Semantisch ähnliche Webseiten, Near-Duplicates und fehlende
Domainwerte werden dadurch jedoch nicht vollständig erfasst. Die zusätzliche Evaluation auf dem unabhängig erhobenen Putra-Datensatz erweitert die
Untersuchung um eine externe Datenquelle. Sie bestätigt den RUT--R0-Vorteil für die
Label-Effizienz mit einer linearen Downstream-Probe, umfasst jedoch nur einen einzelnen
zusätzlichen Webseitenbestand. Die Übertragbarkeit auf weitere Erhebungszeiträume,
Datenquellen und Zielumgebungen bleibt daher offen. Zudem adressiert diese externe
Prüfung die Repräsentations- und Label-Effizienz und nicht das vollständige multimodale
TRI-System oder die URL-first-Kaskade. N5 und N10 variieren die Auswahl der gelabelten Trainingsinstanzen und das
Downstream-Training, während die zugrunde liegenden TAPT-Checkpoints unverändert bleiben.
Die berichteten Unsicherheiten erfassen damit nicht die Variation vollständig neu ausgeführter
selbstüberwachter Vortrainingsläufe.

Die Reduktion von 200.000 auf 50.000 Instanzen betrifft ausschließlich die gelabelte
Downstream-Supervision. Für TAPT werden weiterhin 200.000 Webseiten verarbeitet. Der
PhreshPhish-Bestand stellt grundsätzlich Klasseninformationen bereit; diese werden im
SSL-Pool für das selbstüberwachte Training technisch nicht verwendet und erst für die
späteren Downstream-Experimente wieder zugeordnet. Die Untersuchung quantifiziert damit
den methodischen Nutzen einer ohne Downstream-Labels durchgeführten
Repräsentationsanpassung, nicht den betrieblichen Prozess der erstmaligen Erhebung und
Annotation der Daten. Die verwendete 1-pp-Nichtunterlegenheitsmarge ist eine analytische Entscheidungsschwelle
und keine betrieblich validierte Servicegrenze. Die ergänzende Margensensitivität reduziert
die Abhängigkeit der zentralen FF3-Entscheidung von dieser konkreten Festlegung, ersetzt
jedoch keine anwendungsspezifische Festlegung zulässiger Qualitätsverluste.

Die Klassenbalance der Budgetstichproben wird anhand bereits vorhandener Poollabels
hergestellt. Die Budgetanalyse simuliert daher begrenzte gelabelte Trainingsmengen,
aber keinen vollständig budgetierten Annotationsprozess. Labels für Development,
Kalibration und Evaluation sowie die Information zur klassenweisen Auswahl sind in
den 50.000 Trainingslabels nicht enthalten. Eine betriebliche Einsparung von 75\,\%
des gesamten Annotationsaufwands lässt sich daraus nicht ableiten.

Die fünf internen Testszenarien überlappen; ihre Ergebnisse sind keine fünf
unabhängigen Replikationen. Exakte Sign-Flip-p-Werte beziehen sich auf die vollständige
Enumeration der Vorzeichen, setzen für ihre inferenzielle Interpretation aber eine
entsprechende Symmetrie beziehungsweise Austauschbarkeit der Paardifferenzen voraus.
Mit zehn Replikationen ist diese Annahme nur eingeschränkt prüfbar. Die nachträglichen
Welch-Vergleiche mit XGBoost sind ergänzende Evidenz auf derselben Testpopulation,
keine unabhängige Bestätigung. Ein gespeicherter Analyseplan nach früherem FINAL-Zugriff
ersetzt keine prospektive Vorregistrierung auf einem unberührten Testset.

Die Putra-Erhebungen enden 2023 und liegen damit vor dem verwendeten PhreshPhish-Bestand
aus 2025. Der interne frühe/späte Putra-Schnitt schützt vor direkter Nutzung seiner
Testinstanzen im Zieltraining; die Gesamtauswertung ist jedoch kein prospektiver
Nachweis für zukünftige Angriffe. Eine mögliche Überschneidung mit den ursprünglichen
allgemeinen BERT-/RoBERTa-Vortrainingskorpora wurde nicht überprüft.

Weitere Grenzen betreffen die Eingabeverarbeitung. Der HTML-Textencoder verarbeitet maximal
256 Tokens und erfasst daher nicht notwendigerweise den vollständigen HTML-Inhalt. Der
DOM-Zweig basiert auf einer statischen Rekonstruktion des gespeicherten HTML und nicht auf
einem browsergerenderten Laufzeit-DOM. Die gemessenen Latenz-, Durchsatz- und
Peak-VRAM-Werte gelten ausschließlich für den In-Memory-Modellpfad ohne Netzwerkzugriffe,
Rendering, Datenträger-I/O und vorgelagertes HTML-Parsing.


\section{Betriebliches Einsatzkonzept und Handlungsempfehlung}
\label{sec:betriebskonzept}

\textbf{Einsatzaufgabe und Verantwortlichkeit.}
Konzipiert wird die Vorprüfung verdächtiger Webseiten für die Sicherheitseinheit eines
Kreditinstituts oder dessen IT-Dienstleister. Aus URL, gesichertem HTML und Meldekontext
entsteht ein Klassifikationshinweis mit Modellversion und Schwellenentscheidung für die
fachliche Bearbeitung. Tatsächliche Abläufe, Meldevolumina und Ressourcen der Volksbank
wurden nicht erhoben; der folgende Prozess ist eine eigene Konzeption.

Die technische Vorfilterung großer Ereignismengen für die anschließende menschliche Analyse
wird auch in NIST SP 800-61r3 empfohlen (vgl. [Nelson et al. (2025), S. 25]).
Für die Priorisierung und Validierung sind dort zudem Auswirkungen, Dringlichkeit und
weiterer Fallkontext maßgeblich (vgl. [Nelson et al. (2025), S. 26--27]).
Für den vorliegenden Entwurf folgt daraus eine unterstützende Rolle des Klassifikators:
Die Sicherheitseinheit verantwortet die Fallbewertung und Folgemaßnahmen. Ein hoher
Modellscore allein beschreibt weder den betrieblichen Schaden noch eine kalibrierte
Phishingwahrscheinlichkeit. Automatische Sperrungen sind deshalb kein Bestandteil dieses Konzepts.

Isolierte Erfassung, Netzwerkabruf, Ticketschnittstelle und Prüfoberfläche sind erforderliche
Erweiterungen des Prototyps. Die Eingaben müssen den untersuchten URL-/HTML-/DOM-Daten
entsprechen. Nicht auswertbare Seiten werden gesondert fachlich geklärt. Modellnegative
Seiten erhalten keine Freigabe als sicher. Relevanter Kontext, etwa eine berichtete
Zugangsdatenpreisgabe, kann unabhängig vom Modell eine sofortige Prüfung erfordern.

\begin{figure}[H]
\centering
% Eigene betriebliche Konzeption; KI-unterstuetzt. Keine implementierte Prozessmessung.
\begin{tikzpicture}[
  box/.style={draw=black!55,rounded corners=2pt,align=center,font=\small,
    text width=11.6cm,inner sep=6pt,minimum height=.7cm},
  flow/.style={-{Stealth[length=2mm]},thick},node distance=5mm]
\node[box] (in) {\textbf{Eingang und Aufbereitung}\par
Meldung mit URL und Fallkontext; isolierte Erfassung von HTML/DOM};
\node[box,below=of in,fill=black!7] (model) {\textbf{Gemessener Modellpfad}\par
TRI Full oder URL-first-Kaskade; eingefrorene Modell- und Schwellenversion};
\node[box,below=of model] (review) {\textbf{Prüfsteuerung}\par
Modellalarme und relevante Kontextfälle; Stichprobe modellnegativer Fälle;\par
nicht auswertbare Eingaben separat zur Prüfung};
\node[box,below=of review] (human) {\textbf{Fachliche Entscheidung}\par
Sicherheitseinheit bestätigt Befund und entscheidet über Folgemaßnahmen};
\node[box,below=of human] (learn) {\textbf{Dokumentation und Verbesserung}\par
Geprüfte Labels, Prüfzeit und Fehlerarten; getrennte Entwicklung und Kalibration};
\draw[flow] (in)--node[right,font=\scriptsize]{auswertbar}(model);
\draw[flow] (in.west) -- ++(-.7,0) |- (review.west);
\node[font=\scriptsize,rotate=90] at ($(in.west)!0.5!(review.west)+(-.95,0)$) {nicht auswertbar};
\draw[flow] (model)--(review);
\draw[flow] (review)--(human);
\draw[flow] (human)--(learn);
\end{tikzpicture}

\caption{Konzipierte Einbettung der Webseitenklassifikation in die fachliche Fallprüfung. Nur der grau hinterlegte Modellpfad wurde experimentell untersucht.}
\label{fig:betriebsprozess}
{\footnotesize Eigene Darstellung; KI-unterstützt.}
\end{figure}

\textbf{Modellwahl nach tatsächlichem Engpass.}
Tabelle~\ref{tab:business_model_choice} vergleicht das mit 200.000 Labels trainierte
TRI-System, seine Kaskade und seinen separat gemessenen URL-Zweig.

\begin{table}[H]
\centering
\caption{Gemessene Alternativen und daraus abgeleitete Einsatzentscheidung. TPR/FPR: FINAL; Durchsatz: DEV-In-Memory-Benchmark auf Tesla T4.}
\label{tab:business_model_choice}
\small
\begin{tabular}{lrrrp{5.0cm}}
\toprule
Variante & TPR & FPR & Seiten/s & Konzeptionelle Einordnung \\
\midrule
URL-Pfad & 93,16\,\% & 0,682\,\% & 865,9 & Höchster Durchsatz; mehr übersehene Phishingfälle. \\
TRI Full & 96,12\,\% & 0,666\,\% & 183,2 & Startpunkt bei Qualitätsvorrang und ausreichender Rechenkapazität. \\
TRI-Kaskade & 95,77\,\% & 0,680\,\% & 447,0 & Option bei nachgewiesenem Rechenengpass und akzeptiertem Qualitätsverlust. \\
\bottomrule
\end{tabular}
\end{table}

Bei begrenztem Pilotumfang sprechen höhere TPR und niedrigere FPR für TRI Full.
Die Kaskade erhöht den gemessenen Durchsatz um den Faktor 2,44, erzeugt aber zusätzliche
Fehlalarme und übersehene Phishingfälle. Als Größenordnung entsprechen 10.000 Seiten
beim gemessenen mittleren Durchsatz 54,59 Sekunden für TRI Full und 22,37 Sekunden für
die Kaskade: 32,22 Sekunden Differenz. Diese DEV-Benchmark-Referenz umfasst weder
die gesamte Bearbeitungskette noch eine auf andere Prävalenzen angepasste Laufzeit.
Der Geschwindigkeitsvorteil ist deshalb mit tatsächlichem Eingang, Hardware und Wartezeiten abzuwägen.

\textbf{Prüfmenge und Qualitätsfolgen.}
Die Szenarien aus Abschnitt~\ref{sec:betriebspunkt_szenarien} machen die nachgelagerte
Aufgabe greifbar. Bei 10.000 auswertbaren Seiten und angenommenen 1\,\% Phishing liefert
TRI Full rund 162 Modellalarme: etwa 96 erkannte Phishingfälle und 66 Fehlalarme;
rund vier Phishingfälle bleiben modellnegativ. Die Kaskade erzeugt rechnerisch 1,41
zusätzliche Fehlalarme und 0,35 zusätzliche übersehene Phishingfälle je 10.000 Seiten.
Der schnellere URL-Pfad übersieht bei derselben Annahme 6,84 statt 3,88 Phishingfälle.
Eine kleinere Alarmmenge kann somit auch durch schlechtere Erkennung entstehen.

Werden alle Modellalarme geprüft, beschreibt $A_{\pi}$ zunächst nur diesen Teil der
menschlichen Prüfmenge. Bei einer im Pilot zu messenden mittleren Bearbeitungszeit
$t_A$ in Minuten beträgt der entsprechende Aufwand $A_{\pi}t_A/60$ Stunden.
Kontextfälle, nicht auswertbare Seiten, Negativstichproben und Folgemaßnahmen kommen hinzu.
Eine Personal- oder Kosteneinsparung wurde nicht gemessen. Ebenso wenig sind die rund
660 zusätzlichen Modellpfade der Kaskade im 1-\%-Szenario 660 menschliche Prüffälle.
Der beiliegende Betriebsrechner trennt diese Größen und erlaubt eigene, kenntlich
gemachte Volumen-, Prävalenz- und Zeitannahmen.

\textbf{Datenstrategie bei begrenzter Annotation.}
Die Label-Effizienz liefert eine zweite, von der Laufzeit getrennte Handlungsempfehlung:
Eine vorhandene RUT-Repräsentation ist ein begründeter Ausgangspunkt für begrenztes
gelabeltes Nachtraining. Intern erreicht RUT50k das festgelegte Referenzniveau der
200k-Verfahren; extern erzielt RUT auf Putra über 200 bis 2.000 Ziel-Labels einen
AP-AULC-Vorteil von 2,54 Prozentpunkten. Der interne 50k-Nachweis gilt allerdings für
lineare Probes auf eingefrorenen Repräsentationen. Die in diesem Einsatzkonzept
verwendeten Full-/Kaskadenwerte stammen aus einem mit 200k Labels trainierten System.
Eine Kombination von 50k Trainingslabels mit dessen 96,12\,\% TPR wurde nicht nachgewiesen.

Die 75\,\% Reduktion betrifft zudem ausschließlich Downstream-Trainingslabels;
Entwicklung, Kalibration, Evaluation, Datenerhebung und TAPT bleiben eigene
Aufwandspositionen. Für einen neuen Betrieb sollte zunächst eine Lernkurve mit geprüftem
Zielmaterial erhoben werden. Zusätzliche Ziel-TAPT wäre eine separat zu bewertende Option,
da sie im ergänzenden Putra-Versuch die AP-AULC um 1,14 Prozentpunkte verschlechterte.
Geprüfte Fälle dürfen erst nach Dokumentation und Rollentrennung für spätere Modellstände
genutzt werden. Derselbe Fall darf nicht zugleich Nachtraining und dessen Erfolgskontrolle dienen.

\textbf{Pilotierung und messbare Freigabeentscheidung.}
Als nächster Entwicklungsschritt wird ein Parallelbetrieb vorgeschlagen, bei dem der
bestehende fachliche Prüfprozess die Fallentscheidung weiterhin trifft. TRI Full und
Kaskade bewerten dieselben neu eingehenden auswertbaren Fälle. Modellstand, Schwellen,
Routing und die Erfolgskriterien werden vor der Pilotbewertung festgelegt.
Kalibration auf neuem Zielmaterial und anschließende zeitlich getrennte Bewertung
ermöglichen eine Überprüfung des tatsächlichen Betriebspunkts. Die fachlichen Labels
sollten möglichst ohne Kenntnis der Modellentscheidung vergeben werden.

\begin{table}[H]
\centering
\caption{Vorgeschlagene Prüfkriterien für die betriebliche Weiterentwicklung; noch keine gemessene Freigabe.}
\label{tab:business_pilot}
\small
\begin{tabular}{p{3.0cm}p{10.9cm}}
\toprule
Prüfdimension & Zu erfassender Nachweis und Entscheidung \\
\midrule
Datenabdeckung & Anteil auswertbarer Fälle und Ursachen fehlender Eingaben; gesonderte Bearbeitung der Ausfälle muss möglich sein. \\
Detektionsqualität & TPR, FPR und Precision mit Unsicherheit auf unabhängig geprüften Fällen; die Sicherheitseinheit legt akzeptable Grenzen vor Auswertung fest. Die CAL-Ziel-FPR von 0,5\,\% ist keine bereits bestätigte Betriebsgrenze. \\
Prüfkapazität & Alarme, Kontextfälle, Negativstichprobe, mittlere Prüfzeit und Rückstand; tatsächlicher Aufwand muss in die vereinbarte Kapazität passen. \\
Verarbeitung & Gesamtlatenz einschließlich Erfassung, Durchsatz und Modellrouting; Kaskade nur bei relevantem Ressourcenvorteil und vertretbaren zusätzlichen Fehlern. \\
Änderungskontrolle & Bei Qualitätsverschlechterung zuerst Datenabdeckung und Betriebspunkt prüfen; Rekalibration und Nachtraining getrennt bewerten und an zeitlich späteren Fällen erneut prüfen. \\
\bottomrule
\end{tabular}
\end{table}

Werden nur Modellalarme gelabelt, bleiben übersehene Phishingfälle unbekannt. Der Pilot
benötigt daher entweder eine vollständige fachliche Referenzbewertung oder zusätzlich
eine Zufallsstichprobe modellnegativer Fälle mit dokumentierten Auswahlwahrscheinlichkeiten.
Bei ungleichen Prüfanteilen müssen Prävalenz und Fehlerraten entsprechend gewichtet werden.
Das Konzept verbindet damit die vorhandenen Ergebnisse mit einer überprüfbaren
Einführungsentscheidung: Qualitätsvorrang im ersten Pilot, Kaskade bei nachgewiesenem
Ressourcenbedarf und gezielte Labelverwendung für spätere Modellanpassungen.


\clearpage
\section{Fazit und Ausblick}

Die zusätzliche task-adaptive Anpassung bereits vortrainierter URL- und HTML-Text-Encoder
verbessert im untersuchten Protokoll die Phishing-Erkennung. RUT erreicht bei 20.000 Labels
auf dem vollständigen FINAL-Bestand eine um 6,08 Prozentpunkte höhere mittlere TPR als R0.
Die Adaptationseffekte überlagern sich teilweise.

50.000 Labels bilden die kleinste geprüfte RUT-Stufe, welche die festgelegte
Nichtunterlegenheitsentscheidung gegenüber R0@200k und TF-IDF/XGBoost@200k in allen fünf
Evaluationsbedingungen erfüllt. Die Reduktion um 75\,\% betrifft ausschließlich die gelabelte
Downstream-Trainingsmenge des Linear-Probe-Protokolls; die anschließenden
Full-/Kaskadenwerte beruhen auf 200.000 Trainingslabels. Auf dem externen Putra-Bestand verbessert Quell-RUT die AP-AULC
gegenüber R0 um 2,54 Prozentpunkte; die zusätzliche Putra-Anpassung verschlechtert sie
hingegen um 1,14 Prozentpunkte. Weitere selbstüberwachte Anpassung ist damit kein
allgemeines Verbesserungsversprechen.

HTML-Text, DOM und Gating liefern im sequenziellen DEV-Vergleich jeweils einen zusätzlichen
Beitrag. TRI Full erreicht auf FINAL 96,12\,\% TPR bei 0,666\,\% FPR. Die URL-first-Kaskade
eskaliert 14,93\,\% der Fälle und erreicht 95,77\,\% TPR. Ihr DEV-In-Memory-Durchsatz
ist 2,44-mal so hoch wie der des Vollsystems; die Messung umfasst nicht die gesamte Verarbeitungskette.


Das betriebliche Einsatzkonzept sieht TRI Full zur unterstützenden Webseitenvorprüfung
vor; die Kaskade ist bei nachgewiesenem Rechenengpass zu prüfen. Bei angenommenen
1\,\% Phishing erzeugt TRI Full je 10.000 Seiten rund 162 Modellalarme, darunter 66
Fehlalarme. Modellrouting, Alarmmenge und menschliche Prüfung bleiben getrennte Größen,
deren Übertragung der vorgeschlagene Pilot überprüfen soll.

Feste TAPT-Checkpoints, nachträgliche Analysen und der retrospektive externe Test begrenzen
die Verallgemeinerung. Ausstehend sind unabhängige Vortrainingswiederholungen, zukünftige
externe Systemtests, Kostenmessungen sowie die betriebliche Prüfung von Rekalibration
und Drift-Erkennung.
